\chapter{Digital and analog DNA circuits}

Two inputs are called high, but their concentrations are not exactly one.
One preparation contains $0.91$ units; another contains $1.08$. A downstream
gate must accept both without confusing either with a low input of $0.08$.
Now connect three receivers. Is the upstream signal still high at each one?
A truth table alone cannot answer this question: it has discarded concentration,
time, loading and material inventory.

Chapter 20 supplied reaction syntax and conservation laws. Here we construct a
small circuit interface with explicit quantitative guarantees. Our reference
is deliberately staged: threshold capture finishes before amplification is
evaluated. We derive when this abstraction is composable, then solve a
simultaneous-reaction counterexample that breaks it. The numerical rates are
assigned teaching parameters, not estimates for a particular DNA sequence.

\section{Define a bit as an interval, not a point}

Let $c_\star$ be a chosen concentration scale and write $x=[X]/c_\star$.
All concentrations below are in these normalized units. Define accepted input
sets
\begin{equation}
 \mathcal L=[0,0.1],\qquad \mathcal H=[0.9,1.1].
 \label{eq:bands}
\end{equation}
Values between these sets are \emph{invalid}, not intermediate Boolean values.
Values above $1.1$ also violate the interface: an excessively strong high input
could exhaust a threshold intended to distinguish one high input from two.
The interface therefore has both lower and upper bounds.

A two-input digital contract for a real transfer map $g$ is
\begin{equation}
 g(\mathcal B_i\times\mathcal B_j)\subseteq
 \mathcal B_{b(i,j)},\qquad
 \mathcal B_0=\mathcal L,\quad \mathcal B_1=\mathcal H ,
\end{equation}
where $b$ is the desired Boolean operation. The inclusion must hold for
\emph{every} permitted input and declared parameter perturbation. Checking only
$(0,0),(0,1),(1,0),(1,1)$ does not establish this contract.

For positive weights $w_i$, an integrating stage forms
$s=\sum_iw_ix_i$. The weights are dimensionless conversion factors here;
their physical implementation might require distinct gates, stoichiometry
and inventories. Negative weights do not mean negative molecular concentrations.
One possible signed representation uses two nonnegative rails and decodes their
difference. That representation brings its own cancellation, normalization and
resource obligations.

An OR gate needs its largest low-low sum below its smallest threshold and its
smallest high-containing sum above its largest threshold. With equal weights,
the low-low sum is at most $0.2$, whereas the smallest high-containing sum is
$0.9$. Choose $\theta_{\rm OR}\in[0.25,0.35]$. For AND, the largest incorrect
sum is $1.1+0.1=1.2$, and the smallest correct sum is $0.9+0.9=1.8$.
Choose $\theta_{\rm AND}\in[1.25,1.35]$. These intervals already account for
$0.05$ threshold uncertainty; they are not fitted tolerances.

\section{Capture a finite amount before restoring the signal}

Model threshold capture by $X+Q\to W$, where $Q$ is an input-binding sink.
If this reaction completes before the next stage, its residual input is
\begin{equation}
 r=(s-\theta)_+=\max(s-\theta,0).
 \label{eq:residual}
\end{equation}
This is a material balance: one sink equivalent consumes one input equivalent.
It is not an instantaneous law of molecular recognition.
The drawing in Figure~\ref{fig:capture} shows a minimal complementary-strand
sink, not the full seesaw threshold complex.

\Plate{01-capture}{A finite capture inventory changes free input into a bound
duplex. Arrowheads on continuous backbones mark the $3'$ ends; vertical marks
are base pairing, not covalent bonds. The count invariant is free plus bound
input. The final formula assumes capture has completed and ignores dissociation
over the readout window.}{fig:capture}

Let residual input catalyze conversion of an available output-equivalent
reservoir $F$ into $Y$:
\begin{equation}
 X+F\xrightarrow{k}X+Y,\qquad F\xrightarrow{\lambda}Y.
 \label{eq:restorecrn}
\end{equation}
The second reaction models undesired release. $F$ is a formal reservoir, not
necessarily one physical DNA strand: a molecular implementation must also
account for gate scaffolds, displacement fuel and waste. Here $F+Y=F_0$,
$X=r$ stays constant, and
\begin{align}
 \dot F&=-(kr+\lambda)F,\qquad F(0)=F_0,\\
 Y(t)&=F_0\left(1-e^{-(kr+\lambda)t}\right).
 \label{eq:response}
\end{align}
Time is in seconds; $\lambda$ has units $\mathrm{s}^{-1}$.
For physical bimolecular rate $k_{\rm phys}$, our normalized coefficient is
$k=k_{\rm phys}c_\star$, also in $\mathrm{s}^{-1}$. An increase in output
concentration consumes the reservoir rather than creating signal material.

\Listing{restored}

The implementation uses $\operatorname{expm1}(z)=e^z-1$ to retain small responses
accurately. Input validation rejects negative and nonfinite values. It does not
validate a sequence design or prove the separated stages are a sound chemical
approximation.

Take $F_0=1$, $k=1$, $0\leq\lambda\leq0.001$ and a six-second readout.
The response is monotone increasing in $s$, $\lambda$, $k$, $F_0$ and $t$,
and nonincreasing in $\theta$. Thus extrema over the input and threshold boxes
occur at the corresponding endpoints. Equation~\ref{eq:response}, rather than
a random sampling sweep, certifies the intervals below.

\Listing{gate_bounds}
\begin{center}\input{\Generated/truth.tex}\end{center}

Every computed low output is below $0.1$ and every high output exceeds $0.9$.
All outputs are at most $F_0=1$, so the upper high bound is also preserved.
For the weakest AND input, $r=1.8-1.35=0.45$, giving
$Y(6)=1-e^{-2.7}\simeq0.932794$ even with no helpful leak.

\Plate{02-transfer}{Computed transfer curves for the staged AND reference.
The nominal threshold is $1.3$; the bounding curves include threshold
uncertainty and leakage. Horizontal dashed levels delimit the accepted low
and high regions. This is a finite-time, assigned-rate response, not an
experimental calibration curve.}{fig:transfer}

\section{Derive the readout window}

Waiting longer helps a valid signal release output, but also lets leakage
accumulate. Let accepted output be at most $L$ for low and at least $H$ for
high, with $0<L<H<F_0$. Suppose thresholding leaves zero residual for every low
case and at least $r_{\min}>0$ for every high case. With leak uncertainty
$0\leq\lambda\leq\lambda_{\max}$, conservative constraints are
\begin{align}
 F_0(1-e^{-kr_{\min}t})&\geq H,\\
 F_0(1-e^{-\lambda_{\max}t})&\leq L.
\end{align}
Taking logarithms yields
\begin{equation}
 \frac{-\log(1-H/F_0)}{kr_{\min}}\leq t\leq
 \frac{-\log(1-L/F_0)}{\lambda_{\max}} .
 \label{eq:window}
\end{equation}
For zero maximum leak, the upper bound is infinite in this model. For our AND
reference it is approximately $5.117\leq t\leq105.361$ seconds. Six seconds is
inside the window. If the lower bound exceeds the upper bound, there is
\emph{no certified readout time}; selecting a different threshold in the plot
cannot repair the missing separation without changing the assumptions.

\Listing{read_window}

A real receiver may integrate the entire upstream production history instead
of sampling a free concentration once. Its input is then a trajectory, not just
the terminal scalar in this reference. Gross released signal, free signal,
bound signal and fluorescence are different quantities. The staged contract
must not silently switch between them.

\begin{proposition}[Composition under explicit interfaces]
Consider an acyclic network of staged gates. Suppose each gate maps all valid
input intervals to its required output interval for its assigned read time,
each edge delivers that interval without loading or dilution, and every gate
has the private resources and reset state assumed by its contract. Then every
gate in the network returns the Boolean value of the corresponding ideal
circuit.
\end{proposition}
\begin{proof}
Order gates topologically. The source inputs satisfy their declared intervals.
For the first gate, the transfer inclusion therefore places its output inside
the interval for the correct Boolean value. By the edge assumption the next
gate receives a valid interval. Repeating this argument over the finite order
establishes the claim. No independence assumption about numerical errors is
needed: the inclusion covers all jointly realized valid values.
\end{proof}

The theorem proves the staged model, not an arbitrary one-pot cascade.
Its assumptions carry the physical difficulty. A path of depth $d$ evaluated
in sequential six-second stages has an amplification schedule of $6d$ seconds,
plus capture, transport and measurement costs. This is not the latency of a
continuously reacting DNA circuit. The companion tests exhaust the corners of
the two-layer example $(a\mathbin{\rm OR}b)\mathbin{\rm AND}c$, including both
threshold extremes and both leak extremes.

\section{A kinetic race defeats the ideal threshold}

Remove the scheduling assumption. Let capture and catalytic output release
occur simultaneously:
\begin{equation}
 \dot x=-k_qxq,\quad \dot q=-k_qxq,\quad
 \dot F=-k_cxF,\quad \dot Y=k_cxF .
 \label{eq:race}
\end{equation}
Set $q_0>x_0$ and no leak. An ideal threshold predicts no output, since every
input can eventually be absorbed. But input remains catalytically active
\emph{before} absorption.

\Plate{03-race}{One free input participates in two competing processes.
The capture path removes it; the catalytic path returns it while consuming
output reservoir. Dashed return arrows distinguish catalytic reuse from
creation of extra input. Even when every input is eventually captured, output
already released need not disappear.}{fig:race}

The invariant $q-x=\delta=q_0-x_0>0$ gives
$\dot x=-k_qx(x+\delta)$. Separating variables and using $x(0)=x_0$ gives
\begin{equation}
 x(t)=\frac{\delta x_0e^{-k_q\delta t}}
                 {q_0-x_0e^{-k_q\delta t}}.
\end{equation}
The reservoir depends on the integrated exposure
$A(t)=\int_0^t x(u)\,du$, not merely on the eventual value of $x$. Direct
integration gives
\begin{align}
 A(t)&=\frac1{k_q}
 \log\left(1+\frac{x_0}{\delta}(1-e^{-k_q\delta t})\right),\\
 Y(t)&=F_0(1-e^{-k_cA(t)}),\\
 Y(\infty)&=F_0\left[1-\left(\frac{\delta}{q_0}\right)^{k_c/k_q}\right].
 \label{eq:race-limit}
\end{align}
For $x_0=0.2$, $q_0=0.3$, $F_0=1$, $k_c=k_q=1$, the final output is $2/3$:
a large false positive despite zero leak and complete eventual capture.

\Needspace{16\baselineskip}
\Listing{race_output}
\begin{center}\input{\Generated/race.tex}\end{center}

For the same concentrations, capture ten times faster still gives a limiting
output of about $0.104$, above our low ceiling. A hundredfold rate ratio brings
it below $0.011$. These ratios illustrate a requirement, not a toehold-length
prescription. Sequence, salt, temperature, secondary structure and competing
complexes would determine actual rates. Tests compare the closed form against
an independent four-state Runge--Kutta integration and check both invariants.

\section{Fan-out spends material or changes the load}

A graphical wire can branch without specifying who pays. A molecular signal
cannot. If a fixed inventory $M$ is divided among $n$ equal-volume receiving
compartments, each receives at most $M/n$ molecules. More generally, requiring
concentration $H_i$ in compartment volume $V_i$ demands
\begin{equation}
 M\geq N_A\sum_i H_iV_i
\end{equation}
when $H_i$ is in molar units and $V_i$ in litres. Splitting volume along with
material can preserve concentration, but then each recipient has a smaller
volume: that is a different interface, not free duplication.

\Plate{04-load}{Three equal-volume recipients share a fixed output inventory.
The illustrated partition conserves total material but does not give each
recipient the original concentration. Catalytic buffering can change the
interface only by supplying additional output reservoirs and accounting for
binding occupancy.}{fig:load}

Catalytic reuse does not consume input net, yet a bound catalyst is unavailable
elsewhere. Under a restricted rapid-equilibrium model,
$X+L_i\rightleftharpoons C_i$, with undepleted free loads $\ell_i$ and
dissociation constants $K_i$, we have $[C_i]=x\ell_i/K_i$. Therefore
\begin{equation}
 X_{\rm total}=x+\sum_i[C_i]
 =x\left(1+\sum_i\ell_i/K_i\right),\qquad
 x=\frac{X_{\rm total}}{1+\sum_i\ell_i/K_i}.
 \label{eq:load}
\end{equation}
This approximation requires the loads to remain effectively constant and
binding to equilibrate quickly. It is not a full dynamic receiver model.
Using $X_{\rm total}=0.5$, $\ell_i=K_i=1$, the reduced free catalyst lowers
output at the original six-second readout:
\begin{center}\input{\Generated/loads.tex}\end{center}
The no-load output is valid high; even one stipulated load destroys that
particular guarantee. This effect can occur without off-target hybridization.

\section{Analog computation needs an error model}

Analog computation asks for a numerical value, not membership in a Boolean
interval. For $s>\theta$, the staged transfer has derivative
\begin{equation}
 \frac{\partial Y}{\partial s}
 =F_0kt\,e^{-[k(s-\theta)+\lambda]t}.
 \label{eq:sensitivity}
\end{equation}
At $s=\theta$ the derivative changes discontinuously in the ideal capture
model. Above threshold the response is concave and saturating; it is not a
linear weighted sum. Its global Lipschitz constant with respect to $s$ is at
most $F_0kt$, so small analog errors may be amplified near the threshold even
though valid digital intervals are restored. Digital restoration and analog
accuracy are distinct specifications.

For a cascade with stage error bound $\eta_j$ and Lipschitz bound $L_j$,
let $\epsilon_j$ be the deviation from the ideal numerical trajectory. Then
\begin{equation}
 \epsilon_{j+1}\leq L_j\epsilon_j+\eta_j .
\end{equation}
Repeated substitution exposes products of downstream gains multiplying early
errors. A calibrated low-gain region may be useful for analog arithmetic;
saturated operation may be useful for a digital decision. Neither regime
alone establishes the other. Report the concentration range, calibration
residuals, parameter uncertainty and downstream loading with the result.

\section{Positive feedback is not automatically a memory}

Consider the formal autocatalytic reaction $Y+F\to2Y$ in a closed compartment.
With $N=Y+F$ conserved,
\begin{equation}
 \dot Y=kY(N-Y),\qquad
 Y(t)=\frac{N}{1+(N/Y_0-1)e^{-kNt}}\quad(0<Y_0\leq N).
\end{equation}
The endpoints are equilibria. At zero, the derivative of the vector field is
$kN>0$: a small positive seed grows. At $N$ it is $-kN<0$: the filled state is
attracting. Exact zero remains zero only in the ideal model with no seeding.
Thus there are not two stable states separated by a threshold.

\Listing{feedback}
\Plate{05-feedback}{Computed finite-fuel autocatalysis for two positive seeds
and exact zero. The phase line points toward full conversion for every
interior concentration. Seed-dependent delay is not evidence of bistability;
the dashed zero trajectory is an ideal absorbing state vulnerable to leakage.}{fig:feedback}

For $Y_0<N/2$, the half-conversion time is
$t_{1/2}=\log((N-Y_0)/Y_0)/(kN)$. Smaller seeds delay conversion but do not
change its destination. Genuine bistable memory requires a specified dynamics
with two attracting states, their basins, a write operation and retention
under perturbation. Reuse additionally requires restoring resources and
removing or transforming prior outputs.

Qian and Winfree's seesaw experiments distinguish threshold absorption,
catalytic release, integration and signal restoration~\cite{qian}.
They also distinguish released-signal accounting from instantaneous free
concentration. Our compact transfer model illustrates those engineering
questions, but is not their sequence-level implementation.
More recently, heat-rechargeable circuits use externally supplied thermal
cycles to enable repeated computation~\cite{heat}. That result motivates
explicit reset and energy interfaces; it does not make a closed, depleted
reservoir spontaneously reusable.

\section{Exercises with worked answers}
\begin{enumerate}
\item Why must a high interval have an upper bound? \textbf{Answer:}
A high-plus-low sum must remain below an AND threshold. An arbitrarily large
high input can cross it alone.
\item Find the OR separation inequalities for threshold uncertainty
$\pm0.05$. \textbf{Answer:} $0.2<\theta-0.05$ and
$\theta+0.05<0.9$; any $0.25<\theta<0.85$ strictly separates the sums.
Finite-time amplification further constrains the high margin.
\item Derive the weakest residual for our AND. \textbf{Answer:}
$0.9+0.9-(1.3+0.05)=0.45$.
\item Verify the six-second minimum high output. \textbf{Answer:}
$1-e^{-0.45(6)}=0.932794\ldots$.
\item If maximum leak becomes $0.1$, can the same AND meet both bands?
\textbf{Answer:} The upper time becomes $-\log(0.9)/0.1\simeq1.054$,
below the required $5.117$; the certified window is empty.
\item Explain why $r=0$ after capture does not imply zero concurrent output.
\textbf{Answer:} Output depends on positive exposure $\int x\,dt$ before
capture. Equation~\ref{eq:race-limit} gives a nonzero limit.
\item Compute the race limit for equal capture/release rates and
$(x_0,q_0)=(0.2,0.3)$. \textbf{Answer:} $1-(0.1/0.3)=2/3$.
\item Three equal-volume receivers split $0.95$ normalized output units.
Is each high? \textbf{Answer:} Each gets $0.3167$, outside the high interval.
Preserving concentration requires a different volume or resource allocation.
\item Derive the free fraction with four unit loads and unit $K_i$.
\textbf{Answer:} $x/X_{\rm total}=1/(1+4)=1/5$ under the stated
rapid-equilibrium and undepleted-load assumptions.
\item Is a digital-restoring map globally contractive in analog error?
\textbf{Answer:} Not necessarily. Near the threshold our derivative approaches
$F_0kt=6$, although far above threshold it is small.
\item Does $Y+F\to2Y$ store a robust zero and one? \textbf{Answer:}
No. Zero is unstable to a positive seed; the closed system converts fuel.
Two observed endpoint traces do not prove two stable attractors.
\item Design a gate-composition experiment. \textbf{Rubric:} Measure isolated
and loaded transfer trajectories, calibrate free versus released signal, vary
threshold inventory and input timing, include no-input and off-target controls,
report read windows and uncertainty, and test repeated use separately from
single-use correctness. A truth table at one handpicked time is insufficient.
\end{enumerate}

\section{Vocabulary and the next interface}
\Term{Noise margin}: separation between allowed signal sets and the decision
boundary. \Term{Signal restoration}: mapping every allowed input region back
into a valid output region. \Term{Fan-out}: delivery to multiple downstream
consumers with explicit resource and loading costs. \Term{Retroactivity}:
downstream loading changes upstream dynamics. \Term{Bistability}: two
attracting states under the same fixed external conditions.

Circuit diagrams organize reactions by logical dependency. Chapter 22 instead
asks how spatial bonds and local attachment rules organize computation.
The same discipline survives the change: define the physical interface,
prove the abstract claim, and identify what a molecular realization still
has to demonstrate.
